\documentclass[12pt]{article}
\usepackage[utf8]{inputenc}
\usepackage[T1]{fontenc}
\usepackage{geometry}
\geometry{margin=1in}
\usepackage{enumitem}
\usepackage{hyperref}
\hypersetup{colorlinks=true, linkcolor=blue, urlcolor=blue}

\title{Literature Search Methods Report: Artificial Intelligence in Mental Health Care}
\author{}
\date{}

\begin{document}
\maketitle

\section{Introduction and Research Question}

Artificial intelligence (AI) has rapidly become a transformative force in healthcare, with mental health care representing one of the most promising and complex application domains. Recent systematic reviews have documented AI applications across diagnosis, monitoring, and intervention, including machine learning models for detecting depression, natural language processing for suicide risk prediction, and chatbots for therapeutic support. However, the field remains fragmented across disciplines, and methodological quality varies. To professionalize the search methods for this project, this report documents the literature search conducted to identify studies on AI in mental health care.

The overall project investigates how AI is used in mental health care, with a particular focus on three domains: diagnosis, monitoring, and intervention. The search was designed to be reproducible, transparent, and critically evaluated. It moves beyond the first-year focus on constructing a search strategy by justifying database choices, documenting iterative refinements, and assessing search quality. The following research questions guide the review:

\begin{itemize}
    \item \textbf{Primary research question:} How is AI used in diagnosing mental health illnesses, monitoring disease progression and treatment effectiveness, and conducting AI-assisted mental health interventions?
    \item \textbf{Secondary research question:} What are the limitations, challenges, and ethical concerns in the application of AI technologies in mental health?
\end{itemize}

These questions were formulated to capture both the practical applications of AI in mental health care and the critical considerations surrounding their deployment. The search therefore needed to cover clinical, psychological, and technical literature, and to balance sensitivity with precision across multiple databases.

\section{Databases and Search Platforms}

\section{Search Strategy Development}

\section{Final Search Strategies}

\section{PRISMA-S Flow Diagram}
\section{Critical Evaluation of the Search}

\section{Conclusion}

\section*{Refrences}
% To be completed: Structured reference list of discovered relevant literature.
% Consider using EndNote or Zotero for reference management.

\end{document}
